\documentclass[11pt]{article}

\usepackage[margin=1in]{geometry}
\usepackage[T1]{fontenc}
\usepackage{lmodern}
\usepackage{booktabs}
\usepackage{array}
\usepackage{longtable}
\usepackage{enumitem}
\usepackage{hyperref}
\usepackage{xcolor}
\usepackage{listings}

\hypersetup{
    colorlinks=true,
    linkcolor=blue,
    urlcolor=blue
}

\lstset{
    basicstyle=\ttfamily\small,
    frame=single,
    breaklines=true,
    columns=fullflexible
}

\title{Network Protocols for the Drone Control System}
\author{}
\date{}

\begin{document}

\maketitle

\section{System Overview}

The drone control system contains three main software components:

\begin{enumerate}[leftmargin=*]
    \item The \texttt{current\_frontend} React application running in a browser.
    \item The laptop backend, which acts as a WebSocket-to-UDP bridge.
    \item A Raspberry Pi connected to the laptop through an Ethernet cable.
\end{enumerate}

The frontend communicates with the laptop backend. The backend translates
browser messages into UDP packets for the Raspberry Pi and translates
telemetry packets from the Raspberry Pi back into WebSocket messages for the
browser.

\begin{lstlisting}
Browser frontend
    | WebSocket: ws://127.0.0.1:5173
    v
Laptop backend
    | UDP over Ethernet: Pi port 5000
    v
Raspberry Pi
\end{lstlisting}

\section{Ports and Addresses}

The relevant default settings are:

\begin{center}
\begin{tabular}{>{\ttfamily}l l l}
\toprule
\textnormal{Component or traffic} & \textnormal{Address/port} & \textnormal{Purpose} \\
\midrule
Frontend development server & \texttt{127.0.0.1:6970} & Browser page \\
Frontend to backend & \texttt{ws://127.0.0.1:5173} & WebSocket control connection \\
Laptop to Raspberry Pi & \texttt{192.168.1.50:5000/UDP} & Control and trigger packets \\
Raspberry Pi to laptop & \texttt{<laptop-ip>:5001/UDP} & Telemetry packets \\
\bottomrule
\end{tabular}
\end{center}

The frontend is opened at:

\begin{lstlisting}
http://127.0.0.1:6970
\end{lstlisting}

The backend listens for browser WebSocket connections at:

\begin{lstlisting}
ws://127.0.0.1:5173
\end{lstlisting}

The frontend does not communicate directly with the Raspberry Pi. The
backend performs that translation.

\section{Ethernet Configuration}

When the laptop is connected directly to the Raspberry Pi with an Ethernet
cable, both devices need compatible IP addresses on the same subnet. A
recommended static configuration is:

\begin{center}
\begin{tabular}{l l}
\toprule
Device & IPv4 configuration \\
\midrule
Laptop Ethernet adapter & IP address \texttt{192.168.1.10} \\
Raspberry Pi Ethernet & IP address \texttt{192.168.1.50} \\
Both devices & Subnet mask \texttt{255.255.255.0} \\
Both devices & No gateway required for a direct cable \\
\bottomrule
\end{tabular}
\end{center}

The laptop can test basic reachability with:

\begin{lstlisting}
ping 192.168.1.50
\end{lstlisting}

The existing backend defaults to the Raspberry Pi address
\texttt{192.168.1.50}. If a different address is used, start the backend with
the appropriate value:

\begin{lstlisting}
..\.venv\Scripts\python.exe control_center_backend.py --pi-ip 192.168.1.50
\end{lstlisting}

The backend should print messages similar to:

\begin{lstlisting}
WebSocket: ws://127.0.0.1:5173
Pi UDP target: 192.168.1.50:5000
Telemetry UDP listen: 0.0.0.0:5001
\end{lstlisting}

\section{Control Data Flow}

The frontend sends control messages through its WebSocket connection to the
backend. The backend handles these messages in
\texttt{browser\_handler()}.

For a \texttt{CONTROL} message, the backend:

\begin{enumerate}[leftmargin=*]
    \item Parses the incoming JSON.
    \item Checks that the message version is supported.
    \item Validates roll, pitch, yaw, and throttle values.
    \item Validates the Pico output values.
    \item Adds a sequence number and timestamp.
    \item Serializes the resulting packet as compact UTF-8 JSON.
    \item Sends the bytes to the Raspberry Pi using UDP.
\end{enumerate}

The UDP destination is configured by:

\begin{lstlisting}
--pi-ip           192.168.1.50
--pi-control-port 5000
\end{lstlisting}

The backend uses a datagram socket and sends the packet with
\texttt{sendto()}:

\begin{lstlisting}
payload = json.dumps(message, separators=(",", ":")).encode("utf-8")
self.control_socket.sendto(payload, (self.pi_ip, self.pi_control_port))
\end{lstlisting}

A control packet has the following general form:

\begin{lstlisting}
{
  "version": 1,
  "type": "CONTROL",
  "seq": 42,
  "time_ms": 1728970000000,
  "flight": {
    "roll": 0.0,
    "pitch": 0.2,
    "yaw": 0.0,
    "throttle": 0.5
  },
  "pico": {
    "output0": 0,
    "output1": 0,
    "output2": 0,
    "output3": 0
  }
}
\end{lstlisting}

Trigger commands use the same UDP path. An example is:

\begin{lstlisting}
{
  "version": 1,
  "type": "TRIGGER",
  "name": "valve_1_on",
  "time_ms": 1728970000000
}
\end{lstlisting}

The Raspberry Pi must run a UDP receiver listening on port 5000, conceptually:

\begin{lstlisting}
sock.bind(("0.0.0.0", 5000))
\end{lstlisting}

\section{Telemetry Data Flow}

Communication is bidirectional. The Raspberry Pi can transmit data to the
laptop over the same Ethernet cable.

\begin{lstlisting}
Raspberry Pi
    | UDP to 192.168.1.10:5001
    v
Laptop backend
    | WebSocket broadcast
    v
Browser frontend
\end{lstlisting}

The backend listens on all laptop interfaces at UDP port 5001:

\begin{lstlisting}
self.telemetry_socket.bind(("0.0.0.0", telemetry_port))
\end{lstlisting}

The Raspberry Pi must send telemetry to the laptop's Ethernet IP, not to
\texttt{127.0.0.1}. If the laptop Ethernet address is
\texttt{192.168.1.10}, the destination is:

\begin{lstlisting}
192.168.1.10:5001/UDP
\end{lstlisting}

The backend's telemetry loop:

\begin{enumerate}[leftmargin=*]
    \item Receives a UDP datagram.
    \item Decodes the UTF-8 JSON.
    \item Checks that its type is \texttt{TELEMETRY}.
    \item Adds a bridge-received timestamp.
    \item Broadcasts the telemetry to connected browser clients.
\end{enumerate}

An example telemetry packet is:

\begin{lstlisting}
{
  "version": 1,
  "type": "TELEMETRY",
  "seq": 123,
  "fc": {
    "flight_mode": "STABILIZE",
    "armed": false,
    "battery": {
      "voltage": 11.8,
      "percentage": 82
    },
    "attitude": {
      "roll": 0.01,
      "pitch": -0.02,
      "yaw": 1.57
    }
  }
}
\end{lstlisting}

The laptop firewall must permit inbound UDP traffic on port 5001. Otherwise,
the Raspberry Pi may transmit packets that Windows blocks before the backend
can receive them.

\section{What UDP Is}

UDP, or User Datagram Protocol, is a lightweight transport protocol. It sends
individual datagrams without first establishing a persistent connection.
Each datagram contains the destination address and port and is transmitted
independently.

UDP does not automatically guarantee:

\begin{itemize}[leftmargin=*]
    \item delivery,
    \item ordering,
    \item duplicate prevention,
    \item retransmission, or
    \item confirmation that the receiving application processed the packet.
\end{itemize}

This differs from TCP, which establishes a connection and provides ordered,
reliable delivery by retransmitting missing data.

\section{Why UDP Fits Control and Telemetry}

Flight-control and telemetry streams are continuously updated. The newest
state is normally more useful than an older state that arrives late.

For example:

\begin{lstlisting}
t = 0 ms:    throttle = 0.30
t = 50 ms:   throttle = 0.35
t = 100 ms:  throttle = 0.40
\end{lstlisting}

If the first packet is delayed, retransmitting it after the third packet may
be undesirable. UDP avoids transport-level retransmission of stale data and
has low overhead.

UDP is suitable here because:

\begin{enumerate}[leftmargin=*]
    \item It has low protocol overhead.
    \item It supports frequent small updates.
    \item It works directly over the Ethernet connection.
    \item It is supported by both Windows and Raspberry Pi.
    \item It allows the application to decide how to handle stale packets.
\end{enumerate}

The direct Ethernet cable should make packet loss uncommon, but Ethernet does
not change UDP's delivery guarantees. The application must still implement
appropriate safety behavior.

\subsection{Reliability for One-Time Commands}

One-time triggers, such as valve commands, need more protection than a
continuously repeated control state. If a UDP trigger is lost, the next
control update will not necessarily repeat it.

A robust application-level design is:

\begin{enumerate}[leftmargin=*]
    \item The laptop sends a trigger with a unique command ID.
    \item The Raspberry Pi performs the command and sends an acknowledgement.
    \item The laptop retries if no acknowledgement arrives.
    \item The Raspberry Pi ignores duplicate command IDs.
\end{enumerate}

The current backend sends triggers over UDP but does not yet implement this
acknowledgement and retry protocol.

\section{Comparison with SPI, I2C, and CAN}

UDP is a network transport protocol. SPI, I2C, and CAN are hardware
communication interfaces or buses. They serve different physical and
architectural purposes.

\begin{center}
\begin{tabular}{p{0.18\textwidth} p{0.18\textwidth} p{0.14\textwidth} p{0.18\textwidth} p{0.20\textwidth}}
\toprule
Technology & Communication type & Typical range & Wiring & Typical use \\
\midrule
UDP over Ethernet & Network protocol & Meters to kilometers & Ethernet cable & Laptop-to-Pi communication \\
SPI & Local electrical bus & Centimeters & Clock, data, chip-select, ground & High-speed local peripherals \\
I2C & Local electrical bus & Centimeters & SDA, SCL, ground & Multiple low-speed sensors \\
CAN & Robust differential bus & Tens to hundreds of meters & CAN-H, CAN-L, ground & Vehicle and distributed controllers \\
\bottomrule
\end{tabular}
\end{center}

\subsection{SPI}

SPI normally uses clock, data-in, data-out, chip-select, and ground signals.
It is fast and predictable, but is intended for short board-level
connections. A typical Windows laptop does not expose general-purpose SPI
pins. Therefore, SPI is not a practical direct communication method between
the laptop and Raspberry Pi over an Ethernet cable.

SPI is appropriate for high-speed peripherals directly connected to the
Raspberry Pi, such as certain sensors, ADCs, displays, or local
microcontrollers.

\subsection{I2C}

I2C uses two signal wires, SDA and SCL, and supports multiple addressed
devices on the same local bus. It is useful for devices such as IMUs,
barometers, temperature sensors, EEPROMs, and battery monitors.

I2C is slower and has shorter practical cable lengths than Ethernet. A laptop
would need a special USB-to-I2C adapter to communicate with the Pi over I2C.
It is therefore not a practical replacement for the current laptop-to-Pi
Ethernet link.

\subsection{CAN}

CAN is a robust, message-oriented bus commonly used in vehicles and robotics.
It uses differential signaling, message identifiers, arbitration, and
substantial error detection. CAN is a strong candidate for communication
inside the drone, particularly between the Raspberry Pi and flight
controllers, motor controllers, or battery-management systems.

CAN requires CAN interfaces and transceivers. It is more complex to connect
to a laptop than UDP over Ethernet and generally has lower throughput than
Ethernet. It is best considered a possible internal drone bus rather than a
replacement for the laptop-to-Pi network connection.

\section{Recommended Architecture}

A practical architecture is:

\begin{lstlisting}
Laptop frontend
    | WebSocket
    v
Laptop backend
    | UDP over Ethernet
    v
Raspberry Pi gateway
    | I2C  -> low-speed sensors
    | SPI  -> high-speed local peripherals
    | CAN  -> distributed flight or motor controllers
\end{lstlisting}

The laptop-to-Pi link should remain Ethernet with UDP because the laptop and
Pi already support IP networking. SPI and I2C are better used for local
peripherals attached to the Pi. CAN may be added inside the drone when a
robust multi-controller bus is required.

\section{Safety Requirements}

No communication protocol by itself makes a drone safe. The Raspberry Pi or
flight controller should independently enforce:

\begin{itemize}[leftmargin=*]
    \item a control-packet timeout,
    \item neutral output or motor shutdown when control packets stop,
    \item sequence-number and timestamp validation,
    \item range checks,
    \item duplicate and stale-command handling,
    \item acknowledgement or state confirmation for one-time triggers, and
    \item independent hardware failsafes where possible.
\end{itemize}

The Raspberry Pi should never assume that a successful call to
\texttt{sendto()} proves that the packet was received or acted upon.

\section{Summary}

\begin{itemize}[leftmargin=*]
    \item The frontend runs at \texttt{http://127.0.0.1:6970}.
    \item The frontend connects to the laptop backend at
          \texttt{ws://127.0.0.1:5173}.
    \item The backend sends control and trigger packets to the Pi using UDP
          port 5000.
    \item The Pi sends telemetry to the laptop using UDP port 5001.
    \item A direct Ethernet cable requires compatible IP addresses on both
          devices, such as laptop \texttt{192.168.1.10} and Pi
          \texttt{192.168.1.50}.
    \item UDP is a good fit for frequently changing control and telemetry
          state, but one-time commands need application-level acknowledgements
          and retries.
    \item SPI and I2C are local peripheral buses, while CAN is a robust bus
          that may be useful inside the drone.
\end{itemize}

\end{document}
